\section{Scope of the refined numerical criterion on genuine fields}
\label{sec:s5-nonreal-scope}

\paragraph{Review status.}
The complete post-writing author reading and two independent
whole-proof readings have passed. The exact classical source
interfaces have their own recorded scope.

The strict improvements in MA-RB-013 and MA-RB-014 do not make
their numerical exclusion criterion decisive for every field.
We prove this for the same genuine family constructed in
MA-RB-009. This is a scope result, not another reduction of
the shared-zero gate and not an arithmetic example of
comparison failure.

Keep the three actual representations \(R_j\), conductors
\(N_j\), gamma multiplicities \(d_j^\pm\), and functions
\(H_j,C_j,r_j,h_j,D_j,G_j\) from MA-RB-013/014. Put
\[
 \begin{gathered}
  \lambda_j=\frac12\log N_j,\qquad d_j=d_j^++d_j^-,\qquad
  g_j(u)=d_j^+\psi_{\mathbb R}(u)
                       +d_j^-\psi_{\mathbb R}(u+1),\\
  S_m(u)=\sum_{n=2}^{\infty}\frac{(\log n)^m}{n^u}
                                \qquad(m=1,2,\ u>1).
 \end{gathered}
 \tag{MA-RB.280}
\]
The three pairs \((d_j^+,d_j^-)\) are the fixed pairs
\((4,2),(6,4),(4,6)\). Define the following finite constants,
which therefore do not depend on the number field:
\[
 \begin{aligned}
 A_0&=\frac32\max_j\max\{0,g_j(2)+d_jS_1(2)\},\\
 A_1&=\max_j\max\{0,d_jS_1(3)-g_j(3)\},\\
 B_0&=\max_j\max\left\{0,\frac{3d_j}{4}-g_j(4)
                         +d_j\bigl(4S_2(4)+S_1(4)\bigr)\right\},\\
 A_2&=B_0+\frac{343}{576}A_0,\qquad
 \Lambda_0=\max\left\{A_1,\frac{384}{41}(A_2+1)\right\}.
 \end{aligned}
 \tag{MA-RB.281}
\]
These explicit constants may be bounded by finite sums and
the elementary integral tails of \(S_m\); no zero data
is present in their definitions.

\paragraph{MA-RB-015 (the refined bound tends to \(1/4\)).}
For every actual field \(L\) in the specified Calegari family,
write \(\Delta=\Delta_L\) and
\[
 \lambda_* =\min_j\lambda_j,\qquad
 \mathcal R_L=\frac{\lambda_*-A_1}{10}-1.
 \tag{MA-RB.282}
\]
The actual finite conductors satisfy
\[
                 \log N_j\ge\frac1{120}\log\Delta,
       \qquad \lambda_*\ge\frac1{240}\log\Delta.
 \tag{MA-RB.283}
\]
For each fixed integer \(k\ge1\), if
\(\lambda_*\ge\Lambda_0\) and \(\mathcal R_L\ge2k\),
then all three integer branch tests of MA-RB-014 pass,
and its finite refined envelope satisfies
\[
 \frac14-\frac{24k}{\min_j r_j}
       \le\widehat{\mathfrak K}_{L,k}<\frac14,
 \qquad
 \widehat{\mathfrak K}_{L,k}
       \ge\frac14-\frac{24k}{\mathcal R_L}.
 \tag{MA-RB.284}
\]
In particular, if in addition \(\mathcal R_L\ge192k\),
then \(\widehat{\mathfrak K}_{L,k}\ge1/8\).

Consequently, along any sequence of actual fields in this
family with \(\Delta_L\to\infty\),
\[
          \widehat{\mathfrak K}_{L,k}\longrightarrow\frac14
                                  \quad\hbox{for fixed }k.
 \tag{MA-RB.285}
\]
In particular this applies to the genuine fields \(L_q\)
of (MA-RB.99)--(MA-RB.103). For all sufficiently large
\(q\), neither the test \(r_j<2\) nor the test
\(\widehat{\mathfrak K}_{L_q,1}\le0\) excludes comparison
failure. No assertion of an actual shared zero, failure of
comparison, or positive \(\beta\)-excess follows.

\begin{proof}
\emph{Faithfulness and the actual conductor comparison.}
The checked row (MA-RB.50) gives
\(\chi_W(g)<5\) for every nonidentity element of \(S_5\).
After averaging a positive Hermitian form, all representing
matrices are unitary. A unitary matrix acting as the identity
would have trace five, so \(W\) is faithful. Thus both
\(R_B=\eta\oplus W\) and \(R_T=W\oplus W\eta\)
are faithful. The standard four-dimensional representation
\(V\) is the complement of the constants in the permutation
representation on five letters. If \(V\eta(g)=1\), then
\(V(g)=\eta(g)1\). In the positive-sign case the full
permutation matrix is the identity, so \(g=1\). In the
negative-sign case its trace would be \(1-4=-3\), impossible
for a permutation matrix. Hence \(V\eta\), and therefore
\(R_U=V\eta\oplus X\), are faithful as well.

For every nontrivial subgroup \(I\le S_5\), faithfulness
implies \(\operatorname{codim}R_j^I\ge1\), whereas
\(\operatorname{codim}\operatorname{Reg}^I
   =120(1-|I|^{-1})\le120\).
For \(I=1\) both codimensions vanish. Thus
\[
 \operatorname{codim}R_j^I
       \ge\frac1{120}\operatorname{codim}\operatorname{Reg}^I
                                           \qquad(I\le S_5).
 \tag{MA-RB.286}
\]
Apply this inequality to every actual lower ramification
group in the positive conductor sum (MA-RB.158). It gives
\(a_p(R_j)\ge a_p(\operatorname{Reg})/120\) at each prime,
including the wildly ramified primes. The classical regular
conductor--discriminant identity
\(N(\operatorname{Reg})=\Delta\), with the exact inherited
locator recorded after (MA-RB.168), then proves (MA-RB.283)
after multiplication over primes. No comparison between
an Artin conductor and an unknown automorphic tensor
conductor is used.

\emph{Uniform lower and upper bounds for the real cap.}
Write
\[
             H_j(u)=\lambda_j+g_j(u)+e_j(u),
                         \qquad e_j(u)=\frac{L_j'}{L_j}(u).
 \tag{MA-RB.287}
\]
The actual finite-image coefficient estimate used in
(MA-RB.240) gives, for \(u>1\),
\[
          |e_j(u)|\le d_jS_1(u),\qquad
          |e_j'(u)|\le d_jS_2(u).
 \tag{MA-RB.288}
\]
These are full Euler sums with all ramified factors.
The derivative series is absolutely convergent on every
closed half-line \(u\ge u_0>1\).

The classical trigamma series, DLMF 5.15.1, is positive
on the positive real axis, so each \(g_j\) is increasing.
Also \(S_1\) is decreasing. For \(u\ge3\), therefore,
\[
                      H_j(u)\ge\lambda_j-A_1.
 \tag{MA-RB.289}
\]
For \(1\le u\le3\), the exact paired comparison
(MA-RB.260)--(MA-RB.262), now summed over the whole divisor,
gives
\[
       \frac{H_j(u)}{u-\tfrac12}\ge\frac{H_j(3)}{5/2}.
 \tag{MA-RB.290}
\]
If \(\lambda_j\ge A_1\), the two intervals consequently
give
\[
 \begin{aligned}
 (u-\tfrac12)H_j(u)&\ge\frac{H_j(3)}{10}
                \ge\frac{\lambda_j-A_1}{10}
                                      &&(1\le u\le3),\\
 (u-\tfrac12)H_j(u)&\ge\frac52(\lambda_j-A_1)
                                      &&(u\ge3).
 \end{aligned}
 \tag{MA-RB.291}
\]
Taking the minimum and then using
\(\lfloor C_j\rfloor\ge C_j-1\) proves the lower bound
below. Evaluating that minimum at \(u=2\), and using
(MA-RB.288), proves the upper bound:
\[
       \frac{\lambda_j-A_1}{10}-1
                    \le r_j\le\frac32\lambda_j+A_0.
 \tag{MA-RB.292}
\]
In particular \(\min_j r_j\ge\mathcal R_L\), and the
assumed inequality \(\mathcal R_L\ge2k\) makes all the
integer branch tests pass. Passing these tests is not
a claim that the corresponding zeros exist.

\emph{A uniform derivative sign beyond four.}
The same trigamma series has the elementary upper bound
\[
 0<\psi'(v)=\sum_{n=0}^{\infty}(v+n)^{-2}
                  \le v^{-2}+v^{-1}\qquad(v>0),
 \tag{MA-RB.293}
\]
obtained by bounding the terms after the first with their
decreasing integral. Thus, for \(R\ge4\),
\[
 (R-\tfrac12)g_j'(R)\le\frac{3d_j}{4},
                          \qquad g_j(R)\ge g_j(4).
 \tag{MA-RB.294}
\]
Indeed, \(\psi_{\mathbb R}'(R)=\psi'(R/2)/4
\le1/R^2+1/(2R)\), and the shifted term is no larger
than this bound; multiplying by \(R-\tfrac12\le R\)
proves the first inequality.

For each \(n\ge2\), the function \(R n^{-(R-4)}\)
decreases on \([4,\infty)\), since its logarithmic
derivative is \(1/R-\log n<0\). Therefore
\[
       (R-\tfrac12)S_2(R)\le4S_2(4),
                           \qquad S_1(R)\le S_1(4).
 \tag{MA-RB.295}
\]
Combining (MA-RB.288), (MA-RB.294) and (MA-RB.295)
in the exact derivative of \(h_j\) yields, with
\(x=R-\tfrac12\),
\[
 x^2h_j'(R)\le-\lambda_j+B_0+r_j t(R),\qquad
 t(R)=\frac{(R-\tfrac12)^2(2R-1)}{R^2(R-1)^2}.
 \tag{MA-RB.296}
\]
This estimate retains the sign of the conductor term;
the remaining gamma and Euler terms are bounded uniformly
in the field.

The coefficient \(t\) decreases strictly on \((1,\infty)\).
To check this directly, write
\(t=2x^3/(x^2-1/4)^2\). Then
\[
 \frac{t'}t
       =-\frac{x^2+3/4}{x(x^2-1/4)}<0,
              \qquad t(4)=\frac{343}{576}.
 \tag{MA-RB.297}
\]
Using (MA-RB.292) in (MA-RB.296) consequently gives
\[
 x^2h_j'(R)
    \le-\frac{41}{384}\lambda_j
                    +B_0+\frac{343}{576}A_0
    \le-1\qquad(R\ge4),
 \tag{MA-RB.298}
\]
whenever \(\lambda_j\ge\Lambda_0\).
This is a uniform statement on the whole unbounded interval,
not only the field-dependent eventual derivative sign of
MA-RB-014. Since \(h_j(R)\to0\), strict decrease on
\([4,\infty)\) also implies \(h_j(R)>0\) there.

\emph{Uniform control of the maximizing parameter and denominator.}
For \(1\le\sigma\le4\), every value beyond four is
smaller than \(h_j(4)\), so an actual maximizing point
\(R_*\) in (MA-RB.251) can be chosen with
\[
                    \max\{2,\sigma\}\le R_*\le4.
 \tag{MA-RB.299}
\]
The whole-divisor version of the normalized comparison
(MA-RB.263) gives
\(H_j(\sigma)/(\sigma-\tfrac12)
 \ge H_j(R_*)/(R_*-\tfrac12)\).
Using the exact definition of the maximizing budget therefore
gives
\[
 \begin{aligned}
 G_j(\sigma)
  &=H_j(\sigma)-(\sigma-\tfrac12)
                \frac{H_j(R_*)}{R_*-\tfrac12}
       +\frac{r_j(\sigma-\tfrac12)}{R_*(R_*-1)}\\
  &\ge\frac{r_j(\sigma-\tfrac12)}{12}
                                    \qquad(1\le\sigma\le4).
 \end{aligned}
 \tag{MA-RB.300}
\]
For \(\sigma\ge4\), the strict decrease in (MA-RB.298)
instead makes \(R_* =\sigma\) an actual maximizing point.
The exact formula (MA-RB.267) applies throughout this tail:
\[
 G_j(\sigma)=\frac{r_j(\sigma-\tfrac12)}{\sigma(\sigma-1)},
 \qquad
 \widehat F_{j,k}(\sigma)
    =\frac14+\sigma(\sigma-1)\left(1-\frac{2k}{r_j}\right)
          \ge\frac14\qquad(\sigma\ge4).
 \tag{MA-RB.301}
\]
On \([1,4]\), (MA-RB.300) gives
\[
 \widehat F_{j,k}(\sigma)
       \ge(\sigma-\tfrac12)^2-\frac{24k}{r_j}
       \ge\frac14-\frac{24k}{r_j}.
 \tag{MA-RB.302}
\]
Taking the three infima proves (MA-RB.284), using
\(\min_jr_j\ge\mathcal R_L\); the strict upper bound
is already MA-RB-014. The threshold \(192k\) gives
the stated positive bound \(1/8\).

\emph{The limit on genuine fields and its exact meaning.}
By (MA-RB.283), \(\Delta_L\to\infty\) forces
\(\lambda_*\to\infty\) and \(\mathcal R_L\to\infty\).
For fixed \(k\), all the hypotheses above eventually hold.
The two bounds in (MA-RB.284) squeeze the refined envelope
to \(1/4\), proving (MA-RB.285). More explicitly,
\[
 \mathcal R_L\ge\frac{\log\Delta_L}{2400}
                              -\frac{A_1}{10}-1,
 \tag{MA-RB.303}
\]
so the distance of this envelope from \(1/4\) is
\(O_k(1/\log\Delta_L)\) in this limit.

The family (MA-RB.99)--(MA-RB.103) consists of actual
\(S_5\) normal closures, with the required double
transposition at infinity and unramified prime five,
and \(\Delta_{L_q}\ge q\). Those constructive facts
were proved and independently checked in MA-RB-009;
they are not inferred from a formal zero-order assignment.
Applying the present limit to those fields proves the
asserted nondecisiveness of both tests for all sufficiently
large \(q\). No conclusion about the existence or position
of an actual zero is needed or obtained.
\end{proof}

\paragraph{Classical inputs and finite-cover interpretation.}
DLMF 5.15.1 supplies only the classical trigamma series;
the uniform estimates from it are derived above. The
regular conductor--discriminant identity has the inherited
exact locator Neukirch, VII(11.7)--(11.9), printed
pp.~533--534, already recorded in GRH\_NIGHT NT-C06 and
MA-RB-011/012. No fresh whole-monograph proof audit is
claimed here. The argument applies its positive ramification
filtration formula to the actual faithful finite Artin objects,
not to a proposed analytic zero motive. The finite-cover
conventions remain SGA~1, Expos\'e~V, \S7, and Deligne
(1979), Definition~6.1 and \S\S6.2--6.3, as already fixed
in the preceding sections.

This scope theorem leaves the arithmetic question
MA-RBQ-003 unchanged. The two tests of MA-RB-014 do not
settle all fields; a conclusion for those fields requires
information beyond those tests. No impossibility assertion
about other nonvanishing or comparison methods is made.
